\section{Introduction}

Advances in large language models have proceeded along two axes.
\textbf{Vertical scaling} deepens understanding of a fixed set of high-resource languages
by increasing parameter count, training data, and compute---ChatGPT, Claude, and scaling-law
research belong here.
\textbf{Horizontal scaling} asks instead whether a model can support more languages,
domains, and tasks.
Supporting low-resource tail languages is the canonical horizontal-scaling problem and
is the background of this report.
XLM-R \citep{conneau2020xlmr} is a strong multilingual encoder, but for languages
absent from its training data it suffers from two compounding deficits:
(1)~token over-fragmentation, where a single morpheme is split into many subword pieces,
and (2)~representation mismatch, where the embedding space carries no signal from
the target language.

Glot500 \citep{imani2023glot500} combines vocabulary extension with continued MLM
pretraining to address these deficits for 511 language-scripts.
After extending the tokenizer, XLM-R's embedding matrix $\mathbf{E}\in\mathbb{R}^{250002\times768}$
grows to $\mathbb{R}^{366666\times768}$: the 116,664 new rows have never received a gradient.
Glot500 initializes these rows randomly.
Whether a different starting point---one that reuses information already encoded in
the existing XLM-R embeddings---leads to better convergence and downstream performance
under a limited training budget is the question this report addresses.

We investigate three research questions:
\begin{enumerate}
  \item Does Glot500-style SentencePiece append reduce subword fragmentation for
        XLM-R-unseen tail languages?
  \item With tokenizer and corpus held fixed, does changing only the initialization of
        new token rows alter PPPL and downstream task scores?
  \item Does the sentence representation space exhibit language-identity or
        family/typology structure after continued pretraining?
\end{enumerate}

This experiment is not a full reproduction of Glot500 across 511 languages.
All seven Target7 languages use Latin script, so we make no claim about script diversity.
Instead, we select the smallest XLM-R-unseen tail-language corpus band that still
supports downstream evaluation, and run a controlled 5-way initialization ablation
with identical tokenizer, corpus, objective, schedule, and evaluation protocol.

Our contributions are:
\begin{enumerate}
  \item We propose \famm, a family-aware initialization that places new token rows
        at the centroid of source embeddings belonging to the same language family,
        and verify it against four baselines (\rand, \meanm, \fvt, \wfvt) in a
        controlled 5-way ablation with the tokenizer fixed.
  \item We connect the 2D similarity map's family clustering with downstream results,
        showing that family-aware initialization converges to the best target PPPL,
        Tatoeba retrieval, and NER scores.
  \item We track the step-by-step formation of family geometry in the representation
        space and use it to interpret why family-prior initialization reaches a better
        optimum.
\end{enumerate}
